\documentclass{article}
\usepackage[utf8]{inputenc}
\usepackage[T1]{fontenc}
\usepackage{geometry}
\geometry{
    a4paper,
    left=2.5cm,
    right=2.5cm,
    top=2.5cm,
    bottom=2.5cm
}
\usepackage{graphicx}
\usepackage{float}
\usepackage{hyperref}
\usepackage[english]{babel}
\usepackage{amsmath}
\usepackage{amssymb}
\usepackage{array}
\usepackage{booktabs}
\usepackage{caption}
\usepackage{subcaption}
\usepackage{enumitem}
\usepackage{fancyhdr}
\usepackage{titlesec}
\usepackage{url}
\usepackage{textcomp}
\usepackage{xcolor}
\usepackage{listings}
\usepackage{multirow}
\usepackage{longtable}
\usepackage{tabularx}
\usepackage{times}

% Set up page style
\pagestyle{fancy}
\fancyhf{}
\rfoot{\thepage}
\renewcommand{\headrulewidth}{0pt}

% Title formatting
\titleformat{\chapter}[display]{\normalfont\bfseries}{\chaptertitle\thechapter}{0pt}{\Huge}
\titlespacing*{\chapter}{0pt}{50pt}{40pt}

% Document info
\title{Template Document}
\author{Academic Support Agent}
\date{\today}

\begin{document}
\maketitle

% Logo, Resumen y weas
%\begin{algo}
% Logo, titulo, y cuestiones asi
\twocolumn[
\begin{@twocolumnfalse}
% Logo y cuestiones
\begin{tabular}{cl}
\begin{tabular}{l}
%{ \includegraphics[scale=0.35]{logo-uc-color.jpeg} }
{ \includegraphics[scale=0.28]{logo-uc-byn.png} }
\end{tabular}
& 
\begin{tabular}{l}
\nl
Pontificia Universidad Cat lica de Chile \nl
Facultad de F sica - Departamento de F sica  \nl
FIZ0421 - F\'{i}sica Experimental Avanzada \nl
Profesor: 
\end{tabular}
\end{tabular}

\vspace{0.1cm}

% Ahora el nombre del experimento, los nombres y la fecha
\begin{center}
{\textbf{\huge{Experimento X}}} 
\\
\vspace{0.2cm}
Autores 
\vspace{0.1cm}

%Fecha de entrega
Fecha de Entrega
\end{center}
% aca va el resumen
\begin{abstract}
{\normalsize
Se estudio un plasma ...
} \nl \nl
\end{abstract}		
\end{@twocolumnfalse}]
%\end{algo}

\section{Introducci\'on}
%Aqui va la teoria y las formulas 
\section{Montaje Experimental}
%Dibujos y descripcion de aparataje
\section{Procedimiento}
%Lo que se hizo y como se usaron las formulas
\section{Resultados}
%Los resultados, graficos y tablas
\section{An\'alisis}
%Por que esos resultados y sus errores
\section{Conclusiones}
%Evaluar si se cumplio el objetivo
%Puntos debiles y fuertes (errores por ej)
%Posibles mejoras
\begin{thebibliography}{99}

\end{thebibliography}
% La bibliografia
%\newpage
%\begin{thebibliography}{99}
%\bibitem{ref:LibroChicoLab} David N. Ruzie. (1994). Electric Probes For Low Temperatures Plasmas. New York, New York: The American Vacuum Society.
%\\bibitem{ref:PlasmaEngineering} Reece Roth. (1995). Industrial Plasma Engineering. London: The Institute of Physics.
%\end{thebibliography}
\end{document}